% ══════════════════════════════════════════════════════════════
%  Chapter 0: Quick Start --- See It Run
% ══════════════════════════════════════════════════════════════
\chapter{Quick Start --- See It Run}
\label{ch:quickstart}

\begin{learnbox}
\begin{itemize}
  \item Deploy a multi-node blockchain network with a single command using Docker Compose
  \item Observe blocks being mined, hashed, and propagated between peers in real time
  \item Query the blockchain REST API and view the web admin dashboard
  \item Understand which chapters of the book correspond to each running component
\end{itemize}
\end{learnbox}

\lettrine{T}{ime} required: under 5 minutes. You need only \index{Docker}Docker and \index{Docker Compose}Docker Compose installed.
This chapter exists so you can see the finish line before starting the journey. By the end of these
steps, you will have a multi-node \index{blockchain!network}blockchain network, like the one Bitcoin runs on, running on
your machine---\index{mining}mining blocks, propagating them to peers, and serving a \index{admin dashboard}web admin interface. No Rust
installation required.

% ──────────────────────────────────────────────────────────────
\section{Step 1: Clone the Repository}
\label{sec:step1-clone}
\index{Git!clone}

\begin{minted}[fontsize=\footnotesize]{bash}
git clone https://github.com/bkunyiha/rust-blockchain.git
cd rust-blockchain
\end{minted}

% ──────────────────────────────────────────────────────────────
\section{Step 2: Start the Network}
\label{sec:step2-start}

\begin{minted}[fontsize=\footnotesize]{bash}
cd ci/docker-compose/configs
docker compose up --build
\end{minted}

This builds the \index{Rust!compilation}Rust node from source (first run takes a few minutes while \index{Cargo}Cargo compiles), then
starts two \index{container}containers: a \textbf{\index{miner node}miner} node and a \textbf{\index{webserver}webserver} node.

% ──────────────────────────────────────────────────────────────
\section{Step 3: Watch It Work}
\label{sec:step3-watch}

Once the build finishes, you will see output like this in your terminal:

\begin{minted}[fontsize=\footnotesize,linenos=false]{text}
miner-1     | [INFO] Starting miner on port 8010...
miner-1     | [INFO] Mining block #1 ...
miner-1     | [INFO] Block mined: hash=00a3f2...  height=1  txns=1
webserver-1 | [INFO] Connected to peer miner-1:8010
miner-1     | [INFO] Mining block #2 ...
miner-1     | [INFO] Block mined: hash=0017cb...  height=2  txns=1
webserver-1 | [INFO] Received block #2 from miner-1
\end{minted}

Blocks are being \index{mining}mined, \index{hash}hashed, and propagated across the \index{peer-to-peer!network}network in real time.

% ──────────────────────────────────────────────────────────────
\section{Step 4: Open the Admin Dashboard}
\label{sec:step4-dashboard}

While the network is running, open your browser to:

\begin{minted}[fontsize=\footnotesize,linenos=false]{text}
http://localhost:8080
\end{minted}

If the dashboard prompts you to configure an \index{API key}API key, use the default
\textbf{Admin API key}: \code{admin-secret} (unless you have overridden it via
\code{BITCOIN\_API\_ADMIN\_KEY}).

\begin{notebox}
\textbf{API keys (what they are):} the webserver protects its \index{REST API}REST endpoints
with shared secrets passed on each request. In \index{Docker Compose}Docker Compose there are two keys:

\begin{description}
  \item[Admin key] (\code{BITCOIN\_API\_ADMIN\_KEY}, default \code{admin-secret}):
    full access to \code{/api/admin/*} endpoints used by the admin dashboard.
  \item[Wallet key] (\code{BITCOIN\_API\_WALLET\_KEY}, default \code{wallet-secret}):
    access to wallet-related endpoints.
\end{description}
\end{notebox}

You will see the Web Admin Interface showing the current \index{chain height}chain height, connected \index{peer}peers, recent
\index{block}blocks, and the \index{mempool}mempool. This is the same \index{React}React dashboard we build in Chapter~30.

% ──────────────────────────────────────────────────────────────
\section{Step 5: Query the API}
\label{sec:step5-query}

In another terminal, try:

\begin{minted}[fontsize=\footnotesize]{bash}
curl http://localhost:8080/api/admin/blockchain-info | python3 -m json.tool
\end{minted}

You will get a \index{JSON}JSON response with the \index{chain height}chain height, \index{tip hash}tip hash, peer count, and \index{mempool}mempool size---the
same \index{REST API}API we build in Chapter~24.

% ──────────────────────────────────────────────────────────────
\begin{chaptersummary}
  \item \textbf{Primitives} (Ch~6) --- the \code{\index{block}Block} and \code{\index{transaction}Transaction} structs that were
        serialized into each mined block
  \item \textbf{Cryptography} (Ch~8) --- the \index{SHA-256}SHA-256 hashing and \index{ECDSA}ECDSA signing that produced each
        block hash
  \item \textbf{Chain validation} (Ch~9--10) --- the \index{consensus rules}consensus rules that accepted each block
        into the \index{canonical chain}canonical chain
  \item \textbf{Storage} (Ch~11) --- the \index{sled}sled database that persisted blocks to disk
  \item \textbf{Networking} (Ch~12) --- the \index{TCP}TCP protocol that propagated blocks between the miner
        and webserver nodes
  \item \textbf{Node orchestration} (Ch~13) --- the \index{coordinator}coordinator that routed messages between
        subsystems
  \item \textbf{Web API} (Ch~15) --- the \index{Axum}Axum \index{REST API}REST API that served the \code{/api/admin/*}
        endpoints
  \item \textbf{Web Admin UI} (Ch~21) --- the \index{React}React dashboard you just opened in your browser
  \item \textbf{Docker Compose} (Ch~22) --- the \index{Docker Compose}deployment configuration that wired everything
        together
\end{chaptersummary}

% ──────────────────────────────────────────────────────────────
\section{Clean Up}
\label{sec:cleanup}

\begin{minted}[fontsize=\footnotesize]{bash}
docker compose down -v    # Stop containers and remove volumes
\end{minted}
